\begin{textsample}{NEG Dim 5 – pb – Score -1.00 – pb/pb\_inf\_e\_ef\_1.txt}  \label{ex:f5_neg_195}
Apresentação A Secretaria da Educação do Estado da Paraíba e a União Nacional dos Dirigentes Educacionais ( Undime ) apresentam a Proposta Curricular do Estado da Paraíba para as creches e escolas públicas e privadas das redes estadual e municipal do território paraibano.
A realização deste documento foi decorrente da aprovação da Base Nacional Comum Curricular ( BNCC ) pelo Conselho Nacional de Educação ( CNE ) e a sua homologação, em 2017, pelo Ministério da Educação ( MEC ), que, por sua vez, requereu uma série de iniciativas a ser implementadas pelos entes federados ( União, Estados, Distrito Federal e Municípios ), entre elas, o pacto de colaboração entre Estados e Municípios para se elaborar os Currículos dos Estados, numa perspectiva territorial.
Esse regime de colaboração foi instituído pela Portaria nº 331, de 5 de abril de 2018, que dispõe sobre o Programa de Apoio à Implementação da Base Nacional Comum Curricular - ProBNCC e estabelece diretrizes, parâmetros e critérios para sua implementação.
Assim sendo, para a implantação da BNCC e elaboração dos Currículos Estaduais para os territórios estaduais, houve a participação efetiva e significativa do Ministério da Educação ( MEC ), Conselho Nacional de Secretários de Educação ( Consed ), União Nacional dos Dirigentes Municipais de Educação ( Undime ), Conselho Estadual de Educação ( CEE ) e União Nacional dos Conselhos Municipais de Educação ( UNCME ).
Assim sendo, a Proposta Curricular do Estado da Paraíba foi construída por meio de uma interlocução consistente entre diversidades de vozes e mãos representadas, por professores e educadores das Redes e Sistemas de Ensino do território paraibano, pesquisadores de Universidades Públicas, além de parceiros de movimentos e segmentos sociais.
Logo, é uma proposta curricular legitimada pelo viés democrático, sedimentada no diálogo direto com o professor, que vivencia diariamente experiências com o estudante, que, por sua vez, é sujeito do seu tempo, espaço e cultura local.
Tudo isso explicita o nosso compromisso e a nossa lealdade em garantir os direitos de aprendizagens dos educandos e da melhoria da educação paraibana.
O direito à Educação de qualidade é o centro do planejamento curricular, instrumento norteador das ações escolares nas unidades educacionais públicas ( creches e escolas estaduais e municipais ) e privadas da Paraíba.
Vale salientar que a Proposta Curricular é um documento aberto a ser complementado pelos respectivos Sistemas de Ensino ( público e privado ), por meio de seus currículos, Projeto Político Pedagógico e plano de aula dos professores.
Portanto, esperamos que esta proposta seja uma base de orientação importante para os educadores desenvolverem suas práticas educativas cotidianas, de modo a contribuir para a transposição didática efetiva dos objetos de conhecimento/conteúdos e, consequentemente, para a concretização dos objetivos de aprendizagens traçados, os quais garantirão os direitos de aprendizagem de cada criança, adolescentes e jovens e adultos na Educação Básica no território paraibano.

% matched lemmas: 
\end{textsample}
